\chapter[Marco teórico]{Marco teórico}
\label{cp:marco}
% Extensión estimada: a revisar. El reparto de la ficha daba 27--30 páginas, pero el bloque A cerró en 7
% contra las 10--11 previstas. Ver el encabezado de fichas/02-marco-teorico.md.
% Fuentes: Docs/marco_teorico/Documento_LLMs_VLLMs.tex (recortar NLP general, one-hot, RNN/LSTM a una
% mención breve), Docs/estado_del_arte/resumen_papers.tex (fichas de ~40 trabajos) y
% Docs/tabla_resultados_sota. Los fundamentos clínicos, de calibración y de estadística no están en esos
% documentos.
%
% Feedback del tutor (14/09/2026): el estado del arte pasa a ser la última sección de este capítulo, porque
% revisa y usa los conceptos introducidos antes. Los conjuntos de datos y el sesgo, que antes estaban en el
% estado del arte, son conceptos y por eso van aquí. Se agrega una subsección sobre generación de texto a
% partir de imágenes médicas, que faltaba pese a ser el objeto central del trabajo. La subsección "Uso de un
% VLM generativo como clasificador" se traslada a Materiales y métodos por tratarse del enfoque propio.
%
% Revisión 25/09/2026: la atención pasa a sección propia, antes de visión y lenguaje, porque la usan los
% dos (ficha, decisión 13). El capítulo queda solo con conceptos de computación: los fundamentos clínicos
% se reparten entre Materiales y métodos y las métricas de texto, y los conjuntos concretos pasan a
% Materiales y métodos; aquí quedan el sesgo y la generalización como conceptos (ficha, decisión 14).
% Se agregan los fundamentos del aprendizaje profundo como §2.1 y el capítulo sigue el recorrido del
% pipeline (ficha, «Propuesta del 25/09/2026», aprobada).

Este capítulo reúne los conceptos que usa el resto del trabajo, en el orden en que los atraviesa un caso
dentro del \emph{pipeline} (Figura~\ref{fig:marco_mapa}). Una imagen dermatoscópica llega a un
clasificador, que le asigna una categoría, y a un codificador visual, que recupera casos parecidos; con
esa información, un modelo de lenguaje redacta el reporte. La salida se juzga después con métricas de
clasificación y de texto, y queda por saber si lo medido se sostiene fuera de los datos de entrenamiento.
Los tres componentes son redes neuronales entrenadas con el mismo procedimiento, y casi todos se
construyen sobre la operación de atención; por eso el capítulo empieza por ahí. La última sección revisa
el estado del arte con esos conceptos ya disponibles.

% Figura: mapa del capítulo sobre el pipeline.
% Se usa en la apertura del Capítulo 2. Cada componente lleva la sección que lo explica. Es la versión
% simplificada del diagrama del pipeline de sec:met_pipeline: si uno cambia, revisar el otro.

\begin{figure}[htbp]
    \centering
    \begin{tikzpicture}[x=1mm, y=1mm]

        \node[figentrada, text width=26mm] (img) at (-4, 0) {Imagen\\dermatoscópica};

        \node[figtarea, text width=30mm] (clf) at (40, 13)
            {Clasificador\\{\footnotesize visión, §\ref{sec:marco_vision}}};
        \node[figtarea, text width=30mm] (knn) at (40, -13)
            {Codificador visual y kNN\\{\footnotesize §\ref{sec:marco_vlm}}};

        \node[figtarea, text width=36mm] (gen) at (86, 0)
            {Generador\\{\footnotesize lenguaje, §\ref{sec:marco_llm};\\visión--lenguaje, §\ref{sec:marco_vlm}}};

        \node[figsalida, text width=18mm] (rep) at (126, 0) {Reporte};

        \node[fignodo, text width=30mm, draw=figpurpura!70!black, fill=figpurpura!12] (eval)
            at (126, -24) {Evaluación, §\ref{sec:marco_evaluacion}\\{\footnotesize generalización,
            §\ref{sec:marco_generalizacion}}};

        \draw[figflecha] (img.east) -- ++(5, 0) |- (clf.west);
        \draw[figflecha] (img.east) -- ++(5, 0) |- (knn.west);
        \draw[figflecha] (clf.east) -- node[figetiqueta, above] {categoría} ++(12, 0) |- ([yshift=3mm]gen.west);
        \draw[figflecha] (knn.east) -- node[figetiqueta, below, align=center] {casos\\parecidos} ++(12, 0) |- ([yshift=-3mm]gen.west);
        \draw[figflecha] (gen.east) -- (rep.west);
        \draw[figflecha] (rep.south) -- (eval.north);

        \begin{scope}[on background layer]
            \node[figgrupo={todos los componentes: fundamentos, §\ref{sec:marco_fundamentos};
                  atención, §\ref{sec:marco_atencion}}, fit=(clf) (knn) (gen)] {};
        \end{scope}

    \end{tikzpicture}
    \caption[El capítulo recorrido sobre el \emph{pipeline}]{El capítulo recorrido sobre el
    \emph{pipeline}. Cada componente lleva la sección que presenta los conceptos que usa. El clasificador
    asigna una categoría a la imagen, un codificador visual recupera casos parecidos y el generador redacta
    el reporte con esa información; después la salida se evalúa y se examina si lo medido se sostiene
    fuera de los datos de entrenamiento. Los fundamentos del entrenamiento y la operación de atención son
    comunes a los tres componentes.}
    \label{fig:marco_mapa}
\end{figure}


\section{Fundamentos del aprendizaje profundo}
\label{sec:marco_fundamentos}

Esta sección presenta cómo una red neuronal produce una probabilidad por clase, cómo se ajustan sus
parámetros y cómo se comprueba que lo aprendido sirve fuera de los datos de entrenamiento. Solo incluye
lo que los capítulos siguientes usan sin volver a definirlo.

\subsection{Neurona artificial y perceptrón multicapa}
\label{sec:marco_perceptron}

Una neurona artificial calcula una suma ponderada de sus entradas, le suma un sesgo y aplica a ese
resultado una función no lineal, llamada activación. Una capa reúne varias neuronas que reciben la misma
entrada, y un perceptrón multicapa apila capas de modo que la salida de cada una es la entrada de la
siguiente. La capa $l$ transforma el vector $\mathbf{h}^{(l-1)}$ que recibe en
\begin{equation}
    \mathbf{h}^{(l)} = \phi\!\left(W^{(l)} \mathbf{h}^{(l-1)} + \mathbf{b}^{(l)}\right),
    \label{eq:marco_capa}
\end{equation}
donde $W^{(l)}$ es la matriz de pesos, $\mathbf{b}^{(l)}$ el vector de sesgos y $\phi$ la activación,
aplicada elemento a elemento. Sin la activación, la composición de capas se reduciría a una única
transformación lineal; con ella, la red representa funciones no lineales de la
entrada~\cite{rumelhart1986learning,goodfellow2016deep}.

Los pesos y los sesgos de todas las capas son los parámetros de la red, que se agrupan en $\theta$, y
entrenarla consiste en elegir sus valores. El perceptrón reaparece dentro de arquitecturas más complejas:
la cabeza de clasificación que convierte las características extraídas por una red visual en un puntaje
por clase es una capa de la forma de la Ecuación~\ref{eq:marco_capa}.

\subsection{Clasificación: softmax y entropía cruzada}
\label{sec:marco_softmax}

En un problema de $C$ clases, la última capa devuelve un vector de $C$ puntajes reales sin restricción,
los \emph{logits} $\mathbf{z}$. La función \emph{softmax} los convierte en una distribución de probabilidad,
\begin{equation}
    p_c = \frac{\exp(z_c)}{\sum_{j=1}^{C} \exp(z_j)}, \qquad c = 1, \dots, C,
    \label{eq:marco_softmax}
\end{equation}
con valores positivos que suman uno, y la clase predicha es la de mayor probabilidad. La pérdida de
entropía cruzada mide la probabilidad que el modelo asignó a la clase correcta $y$,
$\mathcal{L} = -\log p_y$. Vale cero cuando toda la probabilidad está en la clase correcta y crece sin
límite a medida que esa probabilidad se acerca a cero~\cite{goodfellow2016deep}.

Varias piezas del trabajo operan sobre estas cantidades. Las modificaciones de la pérdida para el
desbalance de clases (Sección~\ref{sec:marco_transferencia}) cambian la entropía cruzada o los logits
antes de la softmax. La calibración (Sección~\ref{sec:marco_eval_calib}) pregunta si las probabilidades
$p_c$ coinciden con las frecuencias de acierto observadas. Y una de las formas de obtener una categoría de
un modelo generativo, descrita en la Sección~\ref{sec:met_clasif_vlm}, compara las probabilidades que ese
modelo asigna al nombre de cada clase.

\subsection{Entrenamiento por descenso del gradiente}
\label{sec:marco_entrenamiento}

Entrenar la red es buscar los parámetros $\theta$ que minimizan la pérdida promedio sobre el conjunto de
entrenamiento. El descenso del gradiente lo hace de forma iterativa. Calcula el gradiente de la pérdida
respecto de los parámetros y los mueve un paso en la dirección opuesta,
\begin{equation}
    \theta \leftarrow \theta - \eta \, \nabla_{\theta} \mathcal{L},
    \label{eq:marco_gradiente}
\end{equation}
donde $\eta$ es la tasa de aprendizaje. En la variante estocástica, cada paso usa el gradiente estimado
sobre un lote de ejemplos en lugar del conjunto completo, y una época es una pasada por todos los
lotes~\cite{goodfellow2016deep}.

La retropropagación calcula ese gradiente para todos los parámetros a la vez
(Figura~\ref{fig:marco_perceptron}). La pasada hacia adelante evalúa la red capa por capa y guarda los
valores intermedios; la pasada hacia atrás aplica la regla de la cadena desde la pérdida hacia la entrada
y obtiene el gradiente de cada capa a partir del de la capa siguiente. Su costo es del mismo orden que el
de la pasada hacia adelante, y eso hace posible entrenar redes con millones de
parámetros~\cite{rumelhart1986learning,goodfellow2016deep}. A cambio, los valores intermedios y los
gradientes ocupan memoria durante el entrenamiento.

% Figura: perceptrón multicapa, pasada hacia adelante y retropropagación.
% Se usa en sec:marco_fundamentos, subsección de entrenamiento.

\begin{figure}[htbp]
    \centering
    \begin{tikzpicture}[x=1mm, y=1mm,
        neurona/.style = {circle, draw=figtrazo, minimum size=6mm, inner sep=0pt, font=\footnotesize},
    ]

        % Capas.
        \foreach \i in {1, 2, 3} {
            \node[neurona, fill=figazul!15, draw=figazul!80!black] (x\i) at (0, 16 - \i*8) {$x_{\i}$};
        }
        \foreach \i in {1, 2, 3, 4} {
            \node[neurona, fill=fignaranja!18, draw=fignaranja!70!black] (h\i) at (25, 20 - \i*8) {};
        }
        \foreach \i in {1, 2, 3} {
            \node[neurona, fill=figverde!12, draw=figverde!70!black] (z\i) at (50, 16 - \i*8) {$z_{\i}$};
        }

        % Conexiones.
        \foreach \i in {1, 2, 3} { \foreach \j in {1, 2, 3, 4} {
            \draw[figtrazo!60, thin] (x\i) -- (h\j);
        } }
        \foreach \j in {1, 2, 3, 4} { \foreach \k in {1, 2, 3} {
            \draw[figtrazo!60, thin] (h\j) -- (z\k);
        } }

        % Softmax y pérdida.
        \node[figtarea, text width=16mm, minimum height=8mm] (sm) at (72, 0) {softmax};
        \node[figsalida, text width=20mm, minimum height=8mm] (loss) at (100, 0)
            {pérdida\\$\mathcal{L} = -\log p_y$};
        \draw[figflecha] (56, 0) -- (sm.west);
        \draw[figflecha] (sm.east) -- node[figetiqueta, above] {$p_c$} (loss.west);

        % Rótulos de capas.
        \node[font=\footnotesize, text=figtrazo] at (0, -22) {entrada};
        \node[font=\footnotesize, text=figtrazo, align=center] at (25, -22) {capa oculta\\$\phi(W\mathbf{x}+\mathbf{b})$};
        \node[font=\footnotesize, text=figtrazo] at (50, -22) {\emph{logits}};

        % Pasadas.
        \draw[figflecha, figazul!80!black, thick] (0, 24) -- node[figetiqueta, above, text=figazul!80!black]
            {pasada hacia adelante: se calcula la salida} (100, 24);
        \draw[-{Latex[length=2mm, width=1.6mm]}, thick, densely dashed, figbermellon]
            (100, -32) -- node[figetiqueta, below, text=figbermellon]
            {retropropagación: se calcula $\nabla_\theta \mathcal{L}$ capa por capa} (0, -32);

    \end{tikzpicture}
    \caption[Perceptrón multicapa y retropropagación]{Perceptrón multicapa con una capa oculta. La pasada
    hacia adelante transforma la entrada capa por capa hasta los \emph{logits}, que la softmax convierte en
    probabilidades, y la pérdida mide la probabilidad asignada a la clase correcta. La retropropagación
    recorre el camino inverso con la regla de la cadena y obtiene el gradiente de la pérdida respecto de
    cada parámetro.}
    \label{fig:marco_perceptron}
\end{figure}


La tasa de aprendizaje única de la Ecuación~\ref{eq:marco_gradiente} trata igual a todos los
parámetros. Adam adapta el paso a cada uno: mantiene dos promedios móviles por parámetro, uno del
gradiente y otro de su cuadrado, y divide el primero por la raíz del segundo, de modo que los parámetros
con gradientes grandes o ruidosos avanzan con pasos más cortos~\cite{kingma2015adam}. Esos dos estados por
parámetro, sumados al gradiente, son el costo de memoria del ajuste fino que la
Sección~\ref{sec:marco_llm_lora} busca reducir.

\subsection{Generalización y regularización}
\label{sec:marco_regularizacion}

La pérdida de entrenamiento es un medio. Lo que interesa es el desempeño sobre casos que el modelo no
vio, y una red con suficientes parámetros puede memorizar el conjunto de entrenamiento: su pérdida sigue
bajando mientras el error sobre datos nuevos sube. Ese fenómeno se llama
sobreajuste~\cite{goodfellow2016deep}.

Para medirlo, los datos se dividen en tres particiones disjuntas. La de entrenamiento ajusta los
parámetros; la de validación elige entre configuraciones, hiperparámetros y épocas; y la de test se usa
una sola vez, al final, para estimar el desempeño sobre datos nuevos. Una decisión tomada mirando el test
lo convierte en parte del ajuste y vuelve optimista su estimación~\cite{goodfellow2016deep}. Por la misma
razón ningún caso puede aparecer en dos particiones, condición que HAM10000 no cumple en su partición
original (Sección~\ref{sec:met_particiones}).

Las técnicas de regularización reducen el sobreajuste restringiendo lo que la red puede aprender. El
decaimiento de pesos penaliza los parámetros grandes~\cite{goodfellow2016deep}. El \emph{dropout} apaga al
azar una fracción de las neuronas en cada paso de entrenamiento e impide que la red dependa de
combinaciones particulares de ellas~\cite{srivastava2014dropout}. El aumento de datos, que se detalla en
la Sección~\ref{sec:marco_transferencia}, amplía el conjunto con variaciones de las imágenes disponibles.
Conservar los parámetros de la época con mejor desempeño en validación, en lugar de los de la última,
cumple una función parecida~\cite{goodfellow2016deep}.

\subsection{Familias de redes según la estructura de la entrada}
\label{sec:marco_familias}

El perceptrón multicapa trata su entrada como un vector sin estructura: cada neurona se conecta con todas
las entradas y no distingue cuáles son vecinas. Las arquitecturas especializadas incorporan la estructura
de los datos. Las redes neuronales convolucionales (\gls{cnn}) aprovechan que en una imagen los píxeles
cercanos están relacionados y aplican los mismos filtros en toda la imagen~\cite{lecun1998gradient}; se
desarrollan en la Sección~\ref{sec:marco_cnn}. Las redes recurrentes procesan secuencias, como un texto,
un elemento por vez~\cite{elman1990finding}.

Una red recurrente mantiene un estado que resume lo visto hasta el momento y lo actualiza con cada nuevo
elemento, de modo que cada salida depende de toda la secuencia anterior~\cite{elman1990finding}. Para
entrenarla, el gradiente se propaga hacia atrás a lo largo de la secuencia, y en secuencias largas tiende
a desvanecerse, lo que dificulta aprender relaciones entre elementos lejanos. Las redes de memoria a corto
y largo plazo agregan compuertas que controlan qué se conserva en el estado y atenúan ese
problema~\cite{hochreiter1997long}. Aun así, el acceso a una posición lejana sigue pasando por todas las
intermedias, y el cálculo de un paso espera al del anterior, lo que impide paralelizar a lo largo de la
secuencia~\cite{vaswani2017attention}.

\section{Transformers y mecanismo de atención}
\label{sec:marco_atencion}

La atención reemplaza el recorrido paso a paso de las redes recurrentes por una comparación directa, en
la que cada posición de la salida pondera los valores de todas las posiciones de la entrada según la
compatibilidad entre una consulta y una clave. Como todas las posiciones se comparan a la vez, el cálculo
se paraleliza a lo largo de la secuencia. La arquitectura \emph{transformer} se construye sobre
esa operación, y de ella derivan tanto los modelos visuales de la Sección~\ref{sec:marco_vision} como
los modelos de lenguaje de la Sección~\ref{sec:marco_llm}~\cite{vaswani2017attention}.

La atención escalada por producto punto proyecta la entrada en tres espacios, consulta $Q$, clave
$K$ y valor $V$, y devuelve una combinación de valores pesada por la compatibilidad entre consulta
y clave:
\begin{equation}
    \mathrm{Attention}(Q, K, V) = \mathrm{softmax}\!\left(\frac{QK^{\top}}{\sqrt{d_k}}\right) V,
    \label{eq:marco_atencion}
\end{equation}
donde $d_k$ es la dimensión de la clave. La división por su raíz evita que el producto punto crezca
con la dimensión y sature la softmax, que dejaría de discriminar entre
posiciones~\cite{vaswani2017attention}. La Figura~\ref{fig:marco_atencion} muestra la matriz de pesos
que resulta para una secuencia corta.

% Figura: la matriz de pesos de atención.
% Se usa en sec:marco_atencion, después de la Ecuación de la atención escalada por producto punto.
% Los pesos son ilustrativos: muestran la forma de la operación, no un modelo entrenado.

\begin{figure}[htbp]
    \centering
    \begin{tikzpicture}[x=1mm, y=1mm]

        % Consultas (filas) y claves (columnas).
        \foreach \i in {1, 2, 3, 4} {
            \node[fignodo, minimum height=6mm, minimum width=9mm, inner sep=1pt,
                  draw=figazul!80!black, fill=figazul!12] (q\i) at (0, 30 - \i*8) {\footnotesize $q_{\i}$};
            \node[fignodo, minimum height=6mm, minimum width=9mm, inner sep=1pt,
                  draw=fignaranja!70!black, fill=fignaranja!18] (k\i) at (8 + \i*8, 34) {\footnotesize $k_{\i}$};
        }

        % Pesos softmax(QK^T / sqrt(d_k)); cada fila suma uno.
        \foreach \row/\vals in {1/{60,20,10,10}, 2/{15,55,20,10}, 3/{10,25,50,15}, 4/{5,10,25,60}} {
            \foreach \v [count=\col] in \vals {
                \fill[figceleste!\v!white] (4 + \col*8, 26 - \row*8) rectangle (12 + \col*8, 34 - \row*8);
                \node[font=\scriptsize] at (8 + \col*8, 30 - \row*8) {0.\ifnum\v<10 0\fi\v};
            }
        }
        \draw[figtrazo, shift={(12, -6)}] (0, 0) grid[step=8] (32, 32);

        % Salida.
        \foreach \i in {1, 2, 3, 4} {
            \node[fignodo, minimum height=6mm, minimum width=9mm, inner sep=1pt,
                  draw=figverde!70!black, fill=figverde!10, densely dashed] (o\i) at (66, 30 - \i*8)
                  {\footnotesize $o_{\i}$};
            \draw[figflecha] (45, 30 - \i*8) -- (o\i.west);
        }

        \node[font=\footnotesize, text=figtrazo] at (66, 32) {$o_i = \sum_j a_{ij} v_j$};
        \node[font=\footnotesize, text=figtrazo, align=center] at (28, -12)
            {pesos $a_{ij}$ = $\mathrm{softmax}\!\left(q_i k_j^{\top} / \sqrt{d_k}\right)$};

    \end{tikzpicture}
    \caption[Pesos de atención entre las posiciones de una secuencia]{Pesos de atención entre las
    posiciones de una secuencia de cuatro tokens. Cada fila compara la consulta $q_i$ con las claves de
    todas las posiciones, y la softmax convierte esas compatibilidades en pesos que suman uno; la salida
    $o_i$ es la combinación de los valores $v_j$ con esos pesos. Los valores de la matriz son ilustrativos.}
    \label{fig:marco_atencion}
\end{figure}


La atención multi-cabeza proyecta $Q$, $K$ y $V$ en varios subespacios de menor dimensión, aplica
la Ecuación~\ref{eq:marco_atencion} en cada uno por separado y concatena los resultados. Cada cabeza
aprende a capturar relaciones de distinto tipo entre posiciones, y el conjunto amplía la capacidad
de representación sin aumentar la dimensión total del modelo~\cite{vaswani2017attention}.

El transformer original, planteado para traducción automática, combina dos módulos. El codificador
contextualiza la entrada completa, y el decodificador genera la salida token a token atendiendo a lo
codificado y a lo ya generado. Las arquitecturas que usa este trabajo conservan uno solo de los dos:
el ViT de la Sección~\ref{sec:marco_vit} usa el codificador, y los modelos de lenguaje generativos de
la Sección~\ref{sec:marco_llm_grandes}, el
decodificador~\cite{vaswani2017attention,dosovitskiy2021image,brown2020language}.

\section{Aprendizaje profundo para visión}
\label{sec:marco_vision}

El primer componente del \emph{pipeline} recibe la imagen dermatoscópica y devuelve una probabilidad
para cada una de las siete categorías. Para esa tarea el trabajo enfrenta dos familias de arquitectura
de visión. Las redes convolucionales componen
filtros locales apilados y extraen una jerarquía de características. Los transformers visuales
tratan la imagen como una secuencia de tokens y la procesan con la atención de la
Sección~\ref{sec:marco_atencion}. DenseNet y ConvNeXt pertenecen a la primera familia y \gls{vit} y
Swin a la segunda, y las cuatro se comparan en el Capítulo~\ref{cp:res_clasif}.

\subsection{Redes convolucionales}
\label{sec:marco_cnn}

La operación de convolución desliza un filtro de pequeño tamaño sobre la entrada y calcula en cada
posición el producto punto entre el filtro y la ventana que cubre. La conectividad local y el uso de
un mismo filtro en toda la imagen son las dos decisiones de diseño del operador. La primera supone
que la información relevante se agrupa en regiones pequeñas, y la segunda, que un mismo patrón puede
aparecer en cualquier parte (Figura~\ref{fig:marco_convolucion}). Una CNN apila esas operaciones con
activaciones no lineales y reducciones de resolución intermedias~\cite{lecun1998gradient}.

% Figura: la operación de convolución.
% Se usa en sec:marco_cnn, después del primer párrafo.

\begin{figure}[htbp]
    \centering
    \begin{tikzpicture}[x=1mm, y=1mm]

        % Entrada 6x6 con la ventana de 3x3 resaltada.
        \fill[fignaranja!30] (5, 10) rectangle (20, 25);
        \draw[figtrazo!50] (0, 0) grid[step=5] (30, 30);
        \draw[figtrazo, thick] (0, 0) rectangle (30, 30);
        \draw[fignaranja!70!black, very thick] (5, 10) rectangle (20, 25);
        \node[font=\footnotesize, text=figtrazo] at (15, -5) {entrada $6\times 6$};

        % Filtro 3x3.
        \fill[fignaranja!30] (45, 7.5) rectangle (60, 22.5);
        \draw[fignaranja!70!black, shift={(45, 7.5)}] (0, 0) grid[step=5] (15, 15);
        \draw[fignaranja!70!black, very thick] (45, 7.5) rectangle (60, 22.5);
        \node[font=\footnotesize, text=figtrazo, align=center] at (52.5, -5) {filtro $3\times 3$\\(pesos compartidos)};
        \node[font=\large, text=figtrazo] at (37.5, 15) {$\ast$};

        % Salida 4x4 con la celda resultante.
        \fill[figverde!35] (80, 10) rectangle (85, 15);
        \draw[figtrazo!50] (75, 0) grid[step=5] (95, 20);
        \draw[figtrazo, thick] (75, 0) rectangle (95, 20);
        \draw[figverde!70!black, very thick] (80, 10) rectangle (85, 15);
        \node[font=\footnotesize, text=figtrazo, align=center] at (85, -5) {mapa de activación $4\times 4$};
        \node[font=\large, text=figtrazo] at (67.5, 15) {$=$};

        % De la ventana a la celda.
        \draw[figflecha, fignaranja!70!black, bend left=25] (20, 25) to
            node[figetiqueta, above] {producto punto} (80, 15);

    \end{tikzpicture}
    \caption[La operación de convolución]{La operación de convolución. El filtro recorre la entrada y, en
    cada posición, el producto punto entre sus pesos y la ventana que cubre produce una celda del mapa de
    activación. El mismo filtro se aplica en todas las posiciones, de modo que un patrón se detecta
    dondequiera que aparezca.}
    \label{fig:marco_convolucion}
\end{figure}


La profundidad compone esos rasgos locales en rasgos de mayor alcance: los bordes producen texturas,
las texturas producen partes de objetos y las partes producen objetos completos. VGG y ResNet
fijaron el repertorio de bloques que las arquitecturas posteriores refinan, con convoluciones
pequeñas apiladas en el primer caso y conexiones residuales en el segundo, y ese repertorio es el
punto de partida de las redes convolucionales
contemporáneas~\cite{simonyan2015very,he2015residual}.

En DenseNet cada capa recibe las características de todas las anteriores, además de las de la
inmediata, y aporta las suyas a todas las posteriores. Esa reutilización alcanza una capacidad
comparable con menos parámetros que las arquitecturas de conectividad simple, y con conjuntos del
orden de diez mil imágenes, como HAM10000, la economía de parámetros reduce el
sobreajuste~\cite{huang2018densely,he2015residual}.

ConvNeXt moderniza una ResNet adoptando decisiones de diseño de los transformers: convoluciones de
mayor tamaño en lugar de bloques pequeños apilados, menos no-linealidades y normalización por capa.
Con esas modificaciones, una red puramente convolucional iguala al Swin Transformer en exactitud y
escalabilidad. En el comparativo, ConvNeXt representa a las redes convolucionales de la misma época
que el ViT, de modo que el contraste entre familias enfrenta arquitecturas
contemporáneas~\cite{liu2022convnet}.

\subsection{Vision Transformer y Swin Transformer}
\label{sec:marco_vit}

El ViT divide la imagen en una rejilla de parches cuadrados, aplana cada parche y lo proyecta
a un token, y suma a cada token un \emph{embedding} de posición. Sobre esa secuencia aplica un
codificador transformer sin ninguna operación convolucional, y decide la categoría a partir de un
token de clase que se agrega al inicio de la secuencia~\cite{dosovitskiy2021image,vaswani2017attention}.

% Figura 2.1. Del parche al token en el Vision Transformer.
% Se usa en sec:marco_vision, subsección "Vision Transformer y Swin Transformer", después del
% párrafo que define el ViT.

\begin{figure}[htbp]
    \centering
    \begin{tikzpicture}[x=1mm, y=1mm]

        % Imagen dividida en una rejilla de parches.
        \draw[figtrazo, thick, fill=figrelleno] (-26, -12) rectangle (0, 12);
        \foreach \x in {-19.5, -13, -6.5} { \draw[figtrazo] (\x, -12) -- (\x, 12); }
        \foreach \y in {-6, 0, 6} { \draw[figtrazo] (-26, \y) -- (0, \y); }
        \node[font=\small, text=figtrazo, align=center] at (-13, -18) {Imagen dividida\\en parches};

        % Operación de proyección.
        \node[figtarea, text width=30mm] (proj) at (24, 0)
            {Proyectar cada parche\\a un token};

        % Secuencia de tokens.
        \foreach \i in {0, 1, 2, 3, 4} {
            \node[figsalida, minimum width=6mm, minimum height=6mm, inner sep=1pt]
                (t\i) at (52 + \i*9, 6) {\footnotesize $x_{\i}$};
            \node[font=\footnotesize, text=figtrazo] at (52 + \i*9, -3) {$+$};
            \node[font=\footnotesize, text=figtrazo] at (52 + \i*9, -9) {$p_{\i}$};
        }
        \node[font=\small, text=figtrazo, align=center] at (70, -18)
            {Secuencia de tokens con\\embedding de posición por parche};

        % Flechas.
        \draw[figflecha] (0, 0) -- (proj.west);
        \draw[figflecha] (proj.east) -- (46, 0) |- (t0.west);

    \end{tikzpicture}
    \caption[Del parche al token en el Vision Transformer]{Del parche al token en el Vision
    Transformer. La imagen se divide en una rejilla de parches cuadrados y cada parche se aplana y
    proyecta a un token. A cada token se le suma un embedding de posición, porque la operación de
    atención por sí sola no distingue el orden, y la secuencia entra al codificador
    \emph{transformer} sin volver a la retícula de píxeles.}
    \label{fig:marco_vit_patches}
\end{figure}


El ViT no incorpora las dos suposiciones que la convolución da por sentadas, la localidad y la
equivariancia a traslaciones, y tiene que aprenderlas de los datos. Pre-entrenado solo sobre
ImageNet, rinde por debajo de redes convolucionales de tamaño comparable; pre-entrenado sobre
corpus de entre 14 y 300 millones de imágenes, las supera. La variable que decide entre una familia
y otra es la escala del pre-entrenamiento, y por eso importa con qué pesos parte cada arquitectura
antes del ajuste sobre el conjunto médico~\cite{dosovitskiy2021image}.

El Swin Transformer atiende solo dentro de ventanas locales, con costo lineal en el número de
parches en lugar de cuadrático, y desplaza las ventanas entre capas para que la información cruce
de una a otra. Además reduce la resolución por etapas y devuelve una pirámide de mapas de
características, con la estructura jerárquica de una red convolucional y la operación de atención de
un transformer~\cite{liu2021swin}. La Figura~\ref{fig:marco_swin} muestra el desplazamiento de las
ventanas entre dos capas.

% Figura: ventanas y ventanas desplazadas del Swin Transformer. Adaptada de liu2021swin (Figura 2).
% Se usa en sec:marco_vit, en el párrafo del Swin.

\begin{figure}[htbp]
    \centering
    \begin{tikzpicture}[x=1mm, y=1mm]

        % Capa l: cuatro ventanas de 4x4 parches.
        \foreach \c/\x/\y in {figazul/0/16, fignaranja/16/16, figverde/0/0, figpurpura/16/0} {
            \fill[\c!22] (\x, \y) rectangle (\x + 16, \y + 16);
        }
        \draw[figtrazo!40] (0, 0) grid[step=4] (32, 32);
        \draw[figtrazo, very thick] (0, 0) rectangle (32, 32);
        \draw[figtrazo, very thick] (16, 0) -- (16, 32) (0, 16) -- (32, 16);
        \node[font=\footnotesize, text=figtrazo] at (16, -5) {capa $l$: ventanas regulares};

        % Capa l+1: ventanas desplazadas dos parches.
        \begin{scope}[xshift=50mm]
            \foreach \c/\x/\y/\w/\h in {
                figazul/0/24/8/8, fignaranja/8/24/16/8, figverde/24/24/8/8,
                figpurpura/0/8/8/16, figceleste/8/8/16/16, figbermellon/24/8/8/16,
                figazul/0/0/8/8, fignaranja/8/0/16/8, figverde/24/0/8/8} {
                \fill[\c!22] (\x, \y) rectangle (\x + \w, \y + \h);
            }
            \draw[figtrazo!40] (0, 0) grid[step=4] (32, 32);
            \draw[figtrazo, very thick] (0, 0) rectangle (32, 32);
            \draw[figtrazo, very thick] (8, 0) -- (8, 32) (24, 0) -- (24, 32) (0, 8) -- (32, 8) (0, 24) -- (32, 24);
            \node[font=\footnotesize, text=figtrazo] at (16, -5) {capa $l+1$: ventanas desplazadas};
        \end{scope}

        \draw[figflecha] (36, 16) -- (46, 16);

    \end{tikzpicture}
    \caption{Ventanas y ventanas desplazadas del Swin Transformer. Adaptada de \textcite{liu2021swin}.}
    \label{fig:marco_swin}
\end{figure}


\subsection{Transferencia de aprendizaje y desbalance de clases}
\label{sec:marco_transferencia}

Entrenar desde cero una de esas arquitecturas exige más imágenes de las que tiene un conjunto
médico típico. La transferencia de aprendizaje resuelve la escasez aprendiendo primero
representaciones sobre un corpus grande de imágenes naturales y reutilizándolas después para la
tarea de destino. La extracción de características congela el \emph{backbone} y solo entrena la
cabeza de clasificación; el ajuste fino actualiza también los pesos de la red, con una tasa de
aprendizaje menor~\cite{pan2010survey}.

Todas las arquitecturas del trabajo parten de pesos pre-entrenados sobre ImageNet, el corpus de
referencia del reconocimiento visual~\cite{russakovsky2015imagenet}. Con ese punto de partida común,
las diferencias que se miden entre ellas en el Capítulo~\ref{cp:res_clasif} se pueden atribuir a la
arquitectura.

El aumento de datos amplía el corpus de entrenamiento con transformaciones geométricas y
fotométricas de las imágenes disponibles: rotaciones, recortes, desplazamientos y cambios de color.
Funciona como regularizador y reduce el sobreajuste, pero, aplicado por igual a todas las clases,
deja la proporción entre ellas donde estaba~\cite{shorten2019survey}.

La distribución de clases de HAM10000 es muy asimétrica, y un modelo entrenado sin corrección
favorece a las clases mayoritarias~\cite{tschandl2018ham10000}. El Capítulo~\ref{cp:res_clasif}
compara tres modificaciones de la función de pérdida, que actúan durante el entrenamiento y dejan
intacta la métrica de evaluación: ponderar la pérdida por la inversa de la frecuencia de cada clase;
la pérdida focal, que atenúa el aporte de los ejemplos que el modelo ya clasifica con confianza; y
el \emph{balanced softmax}, que suma a cada logit el logaritmo de la frecuencia de su clase en el
entrenamiento~\cite{lin2017focal,ren2020balanced}.

Con esto el clasificador queda completo. Recibe la imagen y devuelve siete probabilidades, ninguna de
las cuales es un texto; el reporte exige un modelo que genere secuencias de palabras, que es el tema de
la sección siguiente.

\section{Modelos de lenguaje}
\label{sec:marco_llm}

Un modelo de lenguaje predice el siguiente token del texto a partir de los anteriores, y de esa
única tarea salen el pre-entrenamiento sobre corpus grandes, el ajuste a una tarea concreta y la
generación de textos nuevos. El generador del \emph{pipeline} pertenece a esta familia. Para el
trabajo interesan cómo se obtiene el modelo, cómo produce el texto y cuánto cuesta ajustarlo.

\subsection{Grandes modelos de lenguaje}
\label{sec:marco_llm_grandes}

Un \gls{llm} se obtiene por pre-entrenamiento auto-supervisado: sobre corpus de miles de
millones de palabras, el modelo aprende a predecir el siguiente token sin anotación humana. De esa
fase sale una distribución de probabilidad sobre el vocabulario condicionada a cualquier prefijo, y
esa distribución es la herramienta que después se reaprovecha para todo lo
demás~\cite{brown2020language,touvron2023llama}.

Dos familias se distinguen por el modo de usar el transformer. La de solo codificador entrena una
representación bidireccional del texto y sirve para clasificar y para comparar; la de solo
decodificador modela la secuencia en una dirección y sirve para generar. El componente de lenguaje
del trabajo pertenece a la segunda~\cite{devlin2019bert,brown2020language,touvron2023llama}.

El comportamiento del modelo se ajusta después en dos etapas. El ajuste fino supervisado lo entrena
sobre pares instrucción--respuesta, y la alineación con \gls{rlhf} lo entrena a preferir las
respuestas que un evaluador humano marca como mejores. Un modelo solo pre-entrenado continúa el
texto que recibe, mientras que uno alineado responde a instrucciones, y el generador del trabajo
recibe su tarea en forma de instrucción~\cite{ouyang2022training}.

La generación recorre la distribución de salida token a token. La búsqueda codiciosa y la búsqueda
por haz eligen las continuaciones más probables y producen textos planos y repetitivos. El muestreo
introduce azar en esa elección, y la temperatura, el \emph{top-k} y el muestreo de núcleo deciden
cuánto. El muestreo de núcleo recorta el vocabulario al subconjunto mínimo que acumula la masa de
probabilidad fijada, y con eso evita la degeneración textual~\cite{holtzman2020curious}. Como la
estrategia de decodificación cambia el texto generado y todas las métricas que se calculen sobre
él, se fija y se documenta como un hiperparámetro más del generador.

\subsection{Ajuste eficiente de parámetros: LoRA}
\label{sec:marco_llm_lora}

El ajuste fino completo de un modelo grande está limitado por la memoria más que por el tiempo de
cómputo. Con el optimizador Adam, cada parámetro entrenado arrastra su gradiente y dos estados
adicionales, y para un modelo de miles de millones de parámetros eso no cabe en la memoria de una
GPU de consumo. Congelar el modelo por completo y entrenar solo una cabeza nueva evita el problema,
pero limita lo que el ajuste puede cambiar~\cite{hu2021lora}.

\gls{lora} ocupa un punto intermedio entre las dos opciones. Sobre los pesos congelados $W$ entrena
un incremento de bajo rango,
\begin{equation}
    W' = W + BA, \qquad B \in \mathbb{R}^{d\times r},\quad A \in \mathbb{R}^{r\times k},\quad
    r \ll \min(d, k),
    \label{eq:marco_lora}
\end{equation}
y solo actualiza $A$ y $B$, que se inicializan de modo que el incremento valga cero al empezar
(Figura~\ref{fig:marco_lora}). El rango $r$ controla la capacidad de adaptación, y valores pequeños
bastan para llevar el modelo a un dominio concreto~\cite{hu2021lora}.

% Figura: reparametrización de LoRA. Adaptada de hu2021lora (Figura 1).
% Se usa en sec:marco_llm_lora, después de la Ecuación de LoRA.

\begin{figure}[htbp]
    \centering
    \begin{tikzpicture}[x=1mm, y=1mm]

        \node[figentrada, text width=10mm, minimum height=7mm] (x) at (0, 0) {$\mathbf{x}$};

        % Rama congelada.
        \node[fignodo, text width=36mm, minimum height=24mm, fill=figrelleno] (W) at (-24, 25)
            {pesos preentrenados\\$W \in \mathbb{R}^{d\times k}$\\[2pt]{\footnotesize congelados}};

        % Rama entrenable: A reduce a r, B vuelve a d.
        \draw[draw=fignaranja!70!black, fill=fignaranja!25, thick] (8, 16) -- (32, 16) -- (26, 22) -- (14, 22) -- cycle;
        \node[font=\footnotesize] at (20, 19) {$A \in \mathbb{R}^{r\times k}$};
        \draw[draw=fignaranja!70!black, fill=fignaranja!25, thick] (14, 28) -- (26, 28) -- (32, 34) -- (8, 34) -- cycle;
        \node[font=\footnotesize] at (20, 31) {$B \in \mathbb{R}^{d\times r}$};
        \node[font=\footnotesize, text=fignaranja!70!black, align=left, anchor=west] at (34, 25)
            {entrenables\\rango $r \ll \min(d,k)$};

        % Suma y salida.
        \node[circle, draw=figtrazo, inner sep=1pt, font=\small] (sum) at (0, 46) {$+$};
        \node[figsalida, text width=26mm, minimum height=7mm] (h) at (0, 58) {$\mathbf{h} = W\mathbf{x} + BA\mathbf{x}$};

        \draw[figflecha] (x.north) -- ++(0, 4) -| (W.south);
        \draw[figflecha] (x.north) -- ++(0, 4) -| (20, 16);
        \draw[figflecha] (20, 22) -- (20, 28);
        \draw[figflecha] (W.north) |- (sum.west);
        \draw[figflecha] (20, 34) |- (sum.east);
        \draw[figflecha] (sum.north) -- (h.south);

    \end{tikzpicture}
    \caption{Reparametrización de LoRA con el incremento de bajo rango $BA$. Adaptada de \textcite{hu2021lora}.}
    \label{fig:marco_lora}
\end{figure}


Los parámetros entrenables y sus estados del optimizador pasan a ser los de $A$ y $B$, una fracción
mínima de los de $W$, y con eso el ajuste vuelve a caber en el presupuesto de una GPU de consumo. Al
terminar, el incremento de la Ecuación~\ref{eq:marco_lora} se suma a $W$ y el modelo resultante se
ejecuta igual que el original, sin ramas ni latencia
adicional~\cite{hu2021lora,yilmaz2025resourceefficient}.

Los modelos generativos que el trabajo ajusta se adaptan con LoRA, y por eso esa parte de la
comparación se hace entera en hardware de consumo. El rango, los módulos congelados y la función de
pérdida del ajuste van en el Capítulo~\ref{cp:metodos}.

Un modelo de lenguaje, sin embargo, solo recibe texto. Para que redacte el reporte a partir de la imagen
hay que conectar su entrada con un codificador visual, y de eso se ocupan los modelos visión--lenguaje.

\section{Modelos visión--lenguaje}
\label{sec:marco_vlm}

\subsection{Alineación contrastiva: CLIP}
\pendiente{Se explica cómo CLIP proyecta imagen y texto a un espacio común, y cómo ese espacio permite la
clasificación \emph{zero-shot} por similitud y el \emph{linear probe}. Es la base de DermLIP.}

\subsection{Modelos multimodales generativos}
\label{sec:marco_vlm_generativos}
\pendiente{Se describe la proyección de características visuales al espacio de un LLM (LLaVA) y el
\emph{visual instruction tuning}, base de LLaVA-Med, Mini-VLM y SkinVL. Se aclara que estos modelos no
reutilizan el esquema codificador--decodificador del transformer original: los tokens visuales proyectados
entran como prefijo a un LLM de solo decodificador.}

\subsection{Generación de texto a partir de imágenes médicas}
\label{sec:marco_vlm_generacion}
\pendiente{Se presenta la tarea central del trabajo: producir texto clínico a partir de una imagen. Se
describen el esquema \emph{image captioning} aplicado al dominio médico y su evolución hacia la redacción de
informes (radiología primero, dermatología después), qué datos requiere (pares imagen--texto clínico), y los
dos problemas característicos de estos modelos: la alucinación de hallazgos y el \emph{hedging}, es decir, la
tendencia a redactar descripciones genéricas y sin compromiso diagnóstico. También se distingue entre la
fidelidad diagnóstica del reporte y su calidad como texto, que retoma las dos condiciones de la
Figura~\ref{fig:intro_dos_tareas} (categoría acertada y descripción fiel a la imagen); sobre esa distinción
se apoya toda la evaluación del Capítulo~\ref{cp:res_gen}.}

\subsection{Generación aumentada por recuperación}
\pendiente{Se explica RAG en general y su variante basada en vecinos visuales (kNN sobre
embeddings).}

\section{Evaluación de modelos}
\label{sec:marco_evaluacion}

\subsection{Métricas de clasificación}
\label{sec:marco_eval_clasif}
\pendiente{Se definen exactitud, F1 ponderado y macro, sensibilidad macro, AUC, y sensibilidad y
especificidad para la tarea binaria, y se muestra por qué la exactitud es engañosa con clases
desbalanceadas.}

\subsection{Calibración de probabilidades}
\label{sec:marco_eval_calib}
\pendiente{Se presentan el diagrama de confiabilidad, ECE, MCE y Brier, y los métodos de calibración post-hoc
\emph{temperature scaling} y \emph{vector scaling}.}

\subsection{Métricas de generación de texto}
\label{sec:marco_eval_texto}
\pendiente{Se describen las métricas con referencia (BLEU, ROUGE-L, METEOR, CIDEr, BERTScore): qué
comparan entre el texto generado y el de referencia, y qué miden y qué no.}

Las métricas con referencia suponen que el texto de referencia es la descripción correcta de la imagen,
y en dermatología ese supuesto es débil. Cuando varios dermatólogos marcan por separado los mismos
criterios dermatoscópicos sobre las mismas imágenes, la concordancia va de pobre a
moderada~\cite{carrera2016validity}. El texto de referencia de un par imagen--texto recoge entonces una
lectura de la lesión entre varias posibles, y una descripción generada puede apartarse de él sin estar
equivocada. Una métrica que compara contra ese único texto la penaliza igual que a una descripción
errónea.

Las métricas sin referencia se apoyan en otra propiedad del dominio. La descripción de una lesión se
escribe con un vocabulario cerrado de términos para color, patrón, forma y distribución, que la
comunidad dermatoscópica estandarizó por consenso. Dentro de ese vocabulario, algoritmos como la lista
de 7 puntos convierten la presencia de ciertos hallazgos en peso diagnóstico, de modo que el camino
entre la imagen y la etiqueta pasa por entidades con
nombre~\cite{kittler2016standardization,argenziano1998checklist}. Con un inventario finito de
entidades se puede contar cuántas de ellas menciona un texto generado sin compararlo con otro texto.

\pendiente{Se definen las dos métricas sin referencia que usa el trabajo, la cobertura de entidades
clínicas y la consistencia diagnóstica, y qué mide y qué no mide cada familia.}

\subsection{Inferencia estadística para comparar modelos}
\pendiente{Se presentan los intervalos por \emph{bootstrap} pareado, la prueba de McNemar y la prueba de
Wilcoxon, y qué pregunta responde cada una. Para Wilcoxon se cita, junto al original, demsar2006statistical,
que justifica su uso para comparar clasificadores.}

\subsection{Interpretabilidad visual}
\pendiente{Se describen los mapas de activación (Grad-CAM) y sus limitaciones.}

\section{Sesgo y generalización}
\label{sec:marco_generalizacion}

\pendiente{Se define el cambio de distribución (\emph{distribution shift}): un modelo aprende la
distribución de su conjunto de entrenamiento, y su desempeño medido sobre una partición de ese mismo
conjunto no garantiza el desempeño sobre imágenes de otra procedencia. Falta la cita del concepto
general (propuesta en refs\_pendientes.md).}

En dermatología el cambio de distribución más visible es el de modalidad. La fotografía clínica
registra la lesión tal como se ve a simple vista, con la piel que la rodea, y la dermatoscopía la
examina con aumento e iluminación controladas, con lo que muestra estructuras que la fotografía no
resuelve, como la red pigmentaria y los patrones
vasculares~\cite{dinnes2018dermoscopy,kittler2016standardization}. Los rasgos visuales tienen
distribuciones distintas en uno y otro registro, y un modelo entrenado sobre una modalidad no se
traslada a la otra sin volver a medirlo. HAM10000 es dermatoscópico y SkinCaRe es de fotografía
clínica, de modo que toda comparación entre los dos conjuntos arrastra esa diferencia de
dominio~\cite{tschandl2018ham10000,shen2025skincare}.

\pendiente{Se presentan las disparidades de desempeño por tono de piel en los conjuntos públicos y la
cuestión de la generalización a poblaciones distintas de la de entrenamiento, que motiva la evaluación
en HIBA (Sección~\ref{sec:met_datasets}).}

\section{Estado del arte}
\label{sec:estado_arte}
\pendiente{Esta sección revisa los trabajos directamente relacionados con la tesina: diagnóstico automático
de lesiones de piel, modelos visión--lenguaje en dermatología y generación de reportes dermatológicos. Cada
bloque cierra indicando qué toma este trabajo de esa línea o en qué se diferencia.}

\subsection{Diagnóstico automático de lesiones de piel}
\pendiente{Un cuadro compara los conjuntos públicos del área (ISIC, PH2, Fitzpatrick17k, SkinCon, Derm1M y
los que usa el trabajo) en modalidad, tamaño, etiquetas, disponibilidad de texto y población; la descripción
detallada de los conjuntos usados está en la Sección~\ref{sec:met_datasets}. Se recorre la evolución desde las CNNs clásicas hacia los transformers, con los resultados de
referencia sobre ISIC y HAM10000 con los que se compara el Capítulo~\ref{cp:res_clasif}, y los problemas
recurrentes del área: desbalance, generalización y sesgo por fototipo.}

\subsection{Modelos visión--lenguaje en dermatología}
\pendiente{Se presentan los modelos pre-entrenados a gran escala con texto clínico (PanDerm, MONET, Derm1M y
DermLIP, MAKE) y su desempeño \emph{zero-shot}, revisión que justifica la elección de DermLIP. Se incluyen
también los VLMs médicos generales efectivamente evaluados en este trabajo o comparables con él (LLaVA-Med y
Phi-3.5-Vision) y la evidencia disponible sobre su desempeño en dermatología.}

\subsection{Generación de reportes dermatológicos}
\pendiente{Se revisan SkinGPT-4, SkinVL (MM-Skin), DermatoLlama/SCALEMED, la destilación de conocimiento de
Chakrabarti et al., Skin-R1, SkinGEN y los ensambles con LLM, junto con la forma en que cada uno evalúa
(métricas con referencia, evaluación por especialistas) y el hardware que requiere. Una tabla comparativa
resume modelo, tarea, datos, métricas y hardware.

La sección cierra con la brecha que aborda esta tesina, que es triple: faltan comparaciones bajo un mismo
protocolo entre clasificadores supervisados, VLMs generales y modelos visión--lenguaje dermatológicos; la
mayoría de los generadores se evalúa por similitud textual y no por fidelidad diagnóstica; y pocos trabajos
consideran la incertidumbre del diagnóstico o el hardware accesible.}
